\chapter{IMPLEMENTATION DETAILS}


\section{Filter Design}

The band-pass filter constitutes the primary analog noise suppression stage of the proposed system. Its function is to pass only the frequencies ranging from 100Hz to 3.4kHz, which is the frequency band for human speech. It removes the frequency components below 100Hz and above 3.4kHz entirely in the analog domain, reducing the computational work on the digital stage. This provides a pre-cleaned signal which can be fed to the TinyML model. Due to the use of analog filtering, near zero latency can be achieved as the filter responds in nanoseconds. 


\subsection{Initial approach}

The filter design started with a first-order passive RC configuration, which was the simplest possible implementation. It required no power supply and only resistors and capacitors were used. Two sections were constructed and cascaded. A high-pass section to block frequencies below 100 Hz, and a low-pass section to block frequencies above 3.4 kHz.

\subsubsection{Component Calculation}

For a passive RC section, the cutoff frequency is given by:
\begin{equation}    
    f_c = \frac{1}{2\pi RC}
\end{equation}
For the high-pass section at \(100\,\text{Hz}\), with \(C = 100\,\text{nF}\) selected as a standard available value, the required resistance is:
\[
R = \frac{1}{2\pi \times 100 \times 100 \times 10^{-9}}
= 15,\!915~\Omega
\]
This was rounded to the nearest standard value:
\[
R = 16~\text{k}\Omega
\]
The actual cutoff frequency with the rounded resistor value is:
\[
f_c = \frac{1}{2\pi \times 16,\!000 \times 100 \times 10^{-9}}
\approx 99~\text{Hz}
\]
For the low-pass section at \(3.4\,\text{kHz}\), with \(C = 10\,\text{nF}\):
\[
R = \frac{1}{2\pi \times 3400 \times 10 \times 10^{-9}}
= 4,\!681~\Omega
\]
This was rounded to the nearest standard value:
\[
R = 4.7~\text{k}\Omega
\]
The actual cutoff frequency with the rounded resistor value is:
\[
f_c = \frac{1}{2\pi \times 4,\!700 \times 10 \times 10^{-9}}
\approx 3,\!386~\text{Hz}
\]
Both cutoff frequencies fall within an acceptable tolerance of the design targets.

\subsubsection{Observed Limitations}

Despite meeting the frequency targets in calculation, the passive RC filter showed some limitations when tested on the breadboard using an oscilloscope and a swept sine wave input.

First, the roll-off rate was insufficient. A single passive RC section produces attenuation of only -20 dB per decade beyond the cutoff frequency. At 50 Hz, which is one octave below the lower cutoff, it was observed that the attenuation was measured at only approximately 6 dB. This is not sufficient to suppress the lower frequencies in the environment. Additionally, noise above the upper cutoff was only attenuated by around 6 dB.

Second, significant passband insertion loss was observed. As a passive filter does not contain an amplifying element, signal energy is lost through the resistors even within the allocated passband. When the high-pass and low-pass sections were cascaded, the total signal loss was measurable and reduced the SNR before the signal could reach the next stage. This is a fundamental limitation of passive RC topologies that cannot be corrected by component selection alone.

Third, the two cascaded sections loaded each other electrically. The output impedance of the high-pass section cannot be neglected in comparison to the input impedance of the low-pass section. When connected in series without buffering, the impedance interaction causes the actual cutoff frequencies to shift from the calculated values. The measured cutoff frequencies on the breadboard deviated from calculated values by more than 15\%, because the single-section formula assumes an ideal voltage source driving the filter, an assumption that fails when two passive sections share the same signal path. These experimental observations confirmed that the passive design was inadequate for the noise environment of the target application and directly motivated the transition to an active filter topology.


\subsection{Active Sallen-Key Band-Pass Filter}

Based on the limitations identified with the passive design, the filter was redesigned using a Sallen-Key active topology. The Sallen-Key configuration was selected for the following reasons specific to this project's constraints.

Firstly, it provides second-order filtering from a single operational amplifier stage, producing a roll-off rate of -40 dB per decade compared to -20 dB per decade from a single passive section. This means that the noise at 50 Hz which is one octave below the 100 Hz cutoff, is attenuated by approximately 14 dB which was a 6 dB attenuation in the passive circuit.

Secondly, the op-amp stage acts as a buffer between the two filter sections. As the output impedance of the op-amp is ideally zero, the high-pass and low-pass stages do not interact through shared impedance. This solves the inter-stage loading problem which was observed in the passive design. In this way, the measured cutoff frequencies match the calculated values.

\subsubsection{Circuit Configuration}

The complete filter consists of two cascaded Sallen-Key stages. The first stage is a second-order high-pass filter which sets the lower cutoff at 100 Hz. The second stage is a second-order low-pass filter setting the upper cutoff at 3.4 kHz. This configuration provides the buffering needed to eliminate inter-stage loading, while also keeping the passband gain at unity (0 dB).

The LM324 operates on a single 5V supply. Because the audio signal is AC-centered around a DC bias point, a virtual ground reference of 2.5 V is established using a resistor voltage divider and stabilised with a bypass capacitor. This mid-supply bias allows the op-amp to faithfully amplify both the positive and negative half-cycles of the audio signal on a single-supply rail.

\subsubsection{High-Pass Stage: Lower Cutoff at 100 Hz}

The high-pass stage uses two equal capacitors (C3 = C4 = 100 nF) and two equal resistors (R3 = R4 = 16 k$\Omega$). The cutoff frequency for the equal-component Sallen-Key high-pass configuration is:
\[
f_c = \frac{1}{2\pi RC}
\]
\[
f_c = \frac{1}{2\pi \times 16,\!000 \times 100 \times 10^{-9}}
\]
\[
f_c = \frac{1}{10.053 \times 10^{-3}}
\]
\[
f_c \approx 99.5~\text{Hz}
\]
This was confirmed in LTspice simulation, where Cursor 2 measured the -6 dB point at 99.808 Hz, in agreement with the calculated value to within 0.3\%.

The choice of C = 100 nF was made because 100 nF is an easily available standard ceramic capacitor value. Given the 100 Hz target, the required resistance works out to approximately 15.9 k$\Omega$. The standard value of 16 k$\Omega$ was selected from the E24 resistor series, as it is the closest available standard value and produces a cutoff less than 0.5\% below the target.

\subsubsection{Low-Pass Stage: Upper Cutoff at 3.4 kHz}

The low-pass stage uses two equal resistors (R1 = R2 = 4.7 k$\Omega$) and two equal capacitors (C1 = C2 = 10 nF). The cutoff frequency is:
\[
f_c = \frac{1}{2\pi RC}
\]
\[
f_c = \frac{1}{2\pi \times 4,\!700 \times 10 \times 10^{-9}}
\]
\[
f_c = \frac{1}{295.3 \times 10^{-6}}
\]
\[
f_c \approx 3,\!386~\text{Hz}
\]
LTspice Cursor 1 confirmed this with a measured -6 dB point at 3.4106 kHz, a deviation of less than 0.7\% from the calculated value.

The capacitor value of 10 nF was selected because it is a standard value that, when paired with a commonly available resistor in the 4-5 k$\Omega$ range, produces the required 3.4 kHz cutoff without requiring any unusual component values. The 4.7 k$\Omega$ resistor is a standard E12 series value available universally.

\subsubsection{Component Summary Table}


\begin{table}[H]
\caption{Calculated and Simulated Cutoff Frequencies}
\label{tab:cutoff_frequencies}
\renewcommand{\arraystretch}{1.4}

\begin{tabularx}{\textwidth}
{|
>{\arraybackslash}L|
>{\arraybackslash}Y|
>{\arraybackslash}Y|
>{\arraybackslash}L|
>{\arraybackslash}L|}

\hline
\textbf{Stage} & \textbf{Components} & \textbf{Values} & \textbf{Calculated $f_c$} & \textbf{Simulated $f_c$} \\
\hline

High-Pass &
$C_3 = C_4,\; R_3 = R_4$ &
$100\,\mathrm{nF},\; 16\,\mathrm{k}\Omega$ &
99.5 Hz &
99.8 Hz \\
\hline

Low-Pass &
$C_1 = C_2,\; R_1 = R_2$ &
$10\,\mathrm{nF},\; 4.7\,\mathrm{k}\Omega$ &
3,386 Hz &
3,411 Hz \\
\hline

\end{tabularx}
\end{table}


\subsubsection{Observed Limitations}

While the second-order Sallen-Key design addressed the fundamental problems of the passive RC filter, when evaluated during simulation and on the breadboard, more limitations were observed.

Firstly, the roll-off rate, which was improved from the passive design, was still found to be insufficient for strong rejection of near-band interference. At -40 dB per decade, noise at 50 Hz is attenuated by approximately 14 dB, and noise at 7 kHz is attenuated by approximately 6 dB. In an industrial environment where fan noise and motor harmonics can sit at 60-80 Hz which is just below the cutoff, the second-order design provides insufficient separation between the noise and the passband edge.

Secondly, the equal-component design yields a quality factor of Q = 0.5 per stage, instead of the Q values needed for a true Butterworth response. The -3 dB point occurs at a frequency higher than the intended cutoff because the two poles of the second-order system are both real and coincident at Q = 0.5. Instead of the -3 dB that a properly positioned Butterworth pole would produce, it was found to be a -6 dB attenuation at the cutoff frequency of 99.8 Hz and 3.41 kHz in the LTspice simulation. This indicates that speech components close to the band edges are attenuated more than intended and that the effective passband is somewhat narrower than intended.

Lastly, the passband flatness in the 2 stage filter, while acceptable, is not maximally flat in the Butterworth sense. The design consisting of equal components allows a small amplitude variation near the cutoff transition which a true pole-placed Butterworth design eliminates. From a hearing assistance perspective, where spectral balance of speech is important for intelligibility, this deviation, which although small, is suboptimal.

\subsection{Fourth-Order Active Sallen-Key Band-Pass Filter}

To address the limitations of the second-order equal-component design, a fourth-order Butterworth band-pass configuration was implemented as four cascaded Sallen-Key stages. The design was developed using LTspice's integrated Filter Wizard tool, which automates the pole placement calculations for higher-order active filters and generates a verified schematic directly. 


\subsubsection{Butterworth Pole Placement}

Two second-order Sallen-Key stages, each with a distinct quality factor Q, are cascaded per side to achieve a fourth-order Butterworth response. Based on the locations of the Butterworth poles on the complex frequency plane, the necessary Q values for a fourth-order Butterworth are:
\[\text{First stage: } Q_1 = 0.5412\]
\[\text{Second stage: } Q_2 = 1.3066\]
These Q values are achieved by using non-equal resistor ratios within each stage rather than the equal-component simplification used in the previous design. The Filter Wizard calculates these ratios by solving the second-order transfer function simultaneously for the target cutoff frequency and the required Q, producing component values that satisfy both constraints exactly. This is the reason the generated schematic contains non-standard resistor values such as 6.04 k$\Omega$, 41.2 k$\Omega$, 147 k$\Omega$, and 174 k$\Omega$, values that would not be arrived at by the simple fc = 1/2$\pi$ RC formula alone.

The result is a true -3 dB point at the cutoff frequencies of 100 Hz and 3.4 kHz, and a maximally flat passband with no amplitude ripple between these frequencies. All speech frequencies from 100 Hz to 3.4 kHz are therefore transmitted with equal amplitude, preventing the filter from introducing any tonal coloration or spectral imbalance into the amplified speech signal.

\subsubsection{Stage Configuration and Component Values}

The filter is divided into four Sallen-Key stages. The high-pass side consists of two stages, referred to as Stage C and Stage D, which together realise the fourth-order high-pass response with a cutoff at 100 Hz. The low-pass side consists of two stages, referred to as Stage A and Stage B, which together realise the fourth-order low-pass response with a cutoff at 3.4 kHz.

Stage C implements the first high-pass Butterworth stage ($Q_1$ = 0.5412) using two equal capacitors of 100 nF and two resistors with values R1 = 6.04 k$\Omega$ and R2 = 41.2 k$\Omega$. The non-equal resistor ratio is what establishes the required Q for this pole pair.

Stage D implements the second high-pass Butterworth stage ($Q_2$ = 1.3066) using two equal capacitors of 10 nF and two resistors with values R1 = 147 k$\Omega$ and R2 = 174 k$\Omega$. The large resistor values in this stage arise from the high Q requirement: achieving $Q_2$ = 1.3066 requires a large ratio between the two resistors, and at 10 nF capacitance the absolute resistor values are pushed into the hundred-K$\Omega$ range. These high resistance values are within the LM324's input bias current specification and do not introduce significant error.

Stage A implements the first low-pass Butterworth stage ($Q_1$ = 0.5412) using capacitors of 100 nF and 10 nF and resistors of R1 = 787 $\Omega$ and R2 = 2.8 k$\Omega$.

Stage B implements the second low-pass Butterworth stage ($Q_2$ = 1.3066) using two equal capacitors of 10 nF and four resistors with values R1 = 3.24 k$\Omega$, R2 = 6.81 k$\Omega$, R3 = 38.3 k$\Omega$, and R4 = 15.8 k$\Omega$.

\subsubsection{Component Summary}

\begin{table}[H]
\caption{Designed Butterworth Filter Stages}
\label{tab:filter_stages}
\renewcommand{\arraystretch}{1.4}

\begin{tabularx}{\textwidth}
{|
>{\arraybackslash}L|
>{\arraybackslash}L|
>{\arraybackslash}Y|
>{\arraybackslash}Y|
>{\arraybackslash}L|
>{\arraybackslash}Y|}

\hline
\textbf{Stage} & \textbf{Function} & \textbf{Capacitors} & \textbf{Resistors} & \textbf{Q} & \textbf{Role} \\
\hline

Stage C &
High-Pass &
$C_1 = C_2 = 100\,\mathrm{nF}$ &
$R_1 = 6.04\,\mathrm{k}\Omega,\; R_2 = 41.2\,\mathrm{k}\Omega$ &
0.5412 &
1st HP Butterworth pole pair \\
\hline

Stage D &
High-Pass &
$C_1 = C_2 = 10\,\mathrm{nF}$ &
$R_1 = 147\,\mathrm{k}\Omega,\; R_2 = 174\,\mathrm{k}\Omega$ &
1.3066 &
2nd HP Butterworth pole pair \\
\hline

Stage A &
Low-Pass &
$C_1 = 100\,\mathrm{nF},\; C_2 = 10\,\mathrm{nF}$ &
$R_1 = 787\,\Omega,\; R_2 = 2.8\,\mathrm{k}\Omega$ &
0.5412 &
1st LP Butterworth pole pair \\
\hline

Stage B &
Low-Pass &
$C_1 = C_2 = 10\,\mathrm{nF}$ &
$R_1 = 3.24\,\mathrm{k}\Omega,\; R_2 = 6.81\,\mathrm{k}\Omega,\; R_3 = 38.3\,\mathrm{k}\Omega,\; R_4 = 15.8\,\mathrm{k}\Omega$ &
1.3066 &
2nd LP Butterworth pole pair \\
\hline

\end{tabularx}
\end{table}

\section{Audio Dataset Preparation}

To develop the TinyML model for transient noise detection, a custom audio dataset was created using recordings captured through commonly available devices such as a mobile phone or laptop. The recordings were recorded as a mono-channel audio at a sampling frequency of 16 kHz. The 16 kHz sampling rate is selected as it satisfies the Nyquist criterion, as the sample frequency 16kHz is more than twice the upper cutoff frequency of the band-pass filter which is 3.4 kHz, which ensures the full speech band is represented without aliasing. Mono channel capture was specified to match the single-channel input configuration of the STM32F401 ADC and to reduce memory and processing requirements during both dataset preparation and model training.


\subsection{Recording Setup}

Recordings were made through a mobile phone or laptop microphone in realistic acoustic conditions rather than a controlled environment. This was a deliberate decision as a model trained on clean studio audio would encounter a systematic mismatch when deployed in factory floors, workshops, and outdoor industrial settings where background noise and room reverberation are always present. Recording under realistic conditions makes the training distribution consistent with the deployment distribution, reducing covariate shift.

Two categories of audio were recorded. The first consists of normal speech and background audio, forming the negative class (label 0). The second consists of recordings containing deliberate transient noise events such as sudden impacts, object drops, door slams, and sharp striking sounds which formed the positive class (label 1). Transient events were produced at varied intensities and distances from the microphone to expose the model to the amplitude variation it will encounter during deployment.


\subsection{Annotation in Audacity}

The recorded audio files were imported into Audacity for manual inspection and annotation. Each recording was visually inspected on the waveform display to identify transient events, which appear as sudden high-amplitude bursts with sharp onset and rapid decay, clearly distinguishable from the lower-amplitude envelope of continuous speech or background noise.

Annotations were created using Audacity's label track feature, which adds a separate annotation layer beneath the waveform without modifying the audio data. For each transient event, a labeled region was drawn across the transient in the waveform view, producing an entry with a start time, end time, and label value. A label value of 1 was assigned to all transient events. Unlabeled regions were treated implicitly as class 0 during frame labeling. The label file format is tab-separated, with each row structured as:

start\_time end\_time label

For example, a transient between 5 and 10 seconds is represented as:

5.000000 10.000000 1

After annotation, the audio was exported as a 16 kHz 16-bit PCM WAV file, and the label track was exported separately using File → Export → Export Labels. Both files share the same base filename for unambiguous matching during processing.


\subsection{Loading and Processing in Python}

The exported WAV files and label files were loaded into Python using soundfile and pandas respectively and operations were performed using numpy. The audio waveform was read as a one-dimensional float array, and the label file was parsed to extract annotated transient regions as start and end time pairs. All time values were converted to sample indices by multiplying by the sampling rate of 16,000, ensuring all subsequent operations were performed in the sample domain to avoid floating-point rounding errors.


\subsection{Frame Segmentation}

Each recording was segmented into fixed-duration frames of 10 ms to match the inference strategy of the embedded STM32F401 system. At 16 kHz, a 10 ms frame contains:
\[
\text{frame\_size} = 0.010 \times 16,\!000 = 160 \text{ samples}
\]
Frames were extracted with 50\% overlap, so each new frame begins 5 ms after the previous one:
\[
\text{hop\_size} = 160 \times 0.5 = 80 \text{ samples}
\]
The overlap ensures transient events near frame boundaries are fully captured by at least one frame. Without overlap, a transient onset falling at a frame boundary could be split such that neither frame contained sufficient energy for reliable classification. The total number of frames for a recording is:
\begin{equation}    
    \text{number\_of\_frames} =
    \left\lfloor
    \frac{\text{total\_samples} - \text{frame\_size}}
    {\text{hop\_size}}
    \right\rfloor + 1
\end{equation}

\subsection{Frame-Level Labeling}

Frame labels were generated by comparing each frame's sample range against the annotated transient regions. A frame spanning indices [frame\_start, frame\_end) is assigned label 1 if any annotated region [region\_start, region\_end) satisfies:
\[
\text{frame\_start} < \text{region\_end} 
\text{ and }
\text{frame\_end} > \text{region\_start}
\]
Otherwise the frame is assigned label 0. This overlapping condition is intentionally permissive as boundary frames containing only the onset or tail of a transient are still labeled 1. This trains the model to trigger the mute circuit slightly before and after each event, which is the safer failure mode: a few extra milliseconds of muted output causes negligible speech loss, whereas passing a transient onset through to the bone conduction actuator defeats the detection system entirely.


% \subsection{Class Imbalance}

% The dataset contains significantly more class 0 frames than class 1 frames. For a 30-minute speech recording with 50\% overlap and 10 ms frames, the total frame count is approximately 360,000. If transient events occupy 30 seconds in total, the class 1 frame count is approximately 6,000, a 60:1 imbalance ratio. Without intervention, a neural network trained on this data learns to predict class 0 for every frame, achieving high apparent accuracy while never detecting a transient.

% To address this, class weights are applied during training using TensorFlow's class\_weight argument in model.fit(). The weight assigned to class 1 is set proportional to the inverse of the class 1 frequency, ensuring each transient frame contributes the same total gradient magnitude as a proportional number of class 0 frames. This forces the model to penalise misclassification of transient events equally to misclassification of speech frames regardless of their relative frequency in the dataset.


\subsection{Export Format}

The processed frames and labels were exported in two formats. The primary format is NumPy binary files (.npy): one file stores the frame data as a two-dimensional float32 array of shape (number\_of\_frames, 160), normalised to [-1.0, 1.0] by dividing by 32,768, and a second file stores the labels as a one-dimensional int8 array of shape (number\_of\_frames,).

The secondary format is a CSV metadata file where each row records the frame index, start time in seconds, end time in seconds, source recording filename, and assigned label. This file allows individual frames to be traced to their source recording for error analysis, and provides a human-readable record of the dataset composition without requiring the full NumPy arrays to be loaded.

The complete pipeline from raw WAV and label files to the NumPy and CSV outputs is implemented as a single reproducible Python script, ensuring the dataset can be regenerated identically if any processing parameter is revised during development.
